\section{TSDB in the context of AirScout}

We ultimately chose TimescaleDB\@. The main reasons were our prior experience
with PostgreSQL and the fact that TimescaleDB provides a powerful and flexible
solution for time-series workloads.

\subsection{How we adapted our project to use TimescaleDB}

To run TimescaleDB in our project, we added a dedicated service to our
\texttt{docker-\linebreak compose} setup, as shown below:

\begin{lstlisting}[style=yaml, caption={Docker-compose service definition for TimescaleDB.}, label={lst:docker-timescaledb-conclusion}]
timescaledb:
  image: timescale/timescaledb:latest-pg14
  container_name: timescaledb
  environment:
    POSTGRES_DB: airscout
    POSTGRES_USER: airscout_user
    POSTGRES_PASSWORD: airscout_pass
  ports:
    - "5434:5432"
  volumes:
    - timescaledb_data:/var/lib/postgresql/data
    - ./AirscoutDB/migrations:/docker-entrypoint-initdb.d/
  networks:
    - airscout-network
  healthcheck:
    test: ["CMD-SHELL", "pg_isready -U airscout_user -d airscout"]
    interval: 10s
    timeout: 5s
    retries: 5
\end{lstlisting}

On its own, this configuration is not sufficient: we also need to enable the
TimescaleDB extension and convert relevant tables into hypertables. A key
benefit for us was that we could initially treat TimescaleDB like a standard
PostgreSQL database. Only after roughly three months of development did we
start using TimescaleDB-specific features such as hypertables and continuous
aggregates. This allowed us to focus on the core functionality first and
optimize later, once the requirements had become clearer.

\subsubsection{How we activated and used TimescaleDB features}\label{sec:timescaledb-features}

To enable the TimescaleDB extension and create hypertables, we added a
dedicated migration.

We enabled the extension using the following SQL statement:

\begin{lstlisting}[language=PGTSSQL, caption={SQL command to activate the TimescaleDB extension.}, label={lst:sql-activate-timescaledb}]
CREATE EXTENSION IF NOT EXISTS timescaledb CASCADE;
\end{lstlisting}

This creates the TimescaleDB extension if it is not already installed. The
\texttt{CASCADE} option ensures that any required dependencies are created as
well. The same approach can be used to enable additional PostgreSQL extensions,
for example PostGIS for geospatial features.

With the extension enabled, the database can use TimescaleDB functionality, but
existing tables do not automatically benefit from performance improvements. To
take advantage of TimescaleDB optimizations, we need to convert suitable tables
into hypertables~\cite{TimescaleDB:02}.

To convert the \texttt{meas} table, we first had to adjust the primary key.
Because we wanted to partition by time, a primary key on \texttt{id} alone was
not ideal. Instead, we used a composite primary key consisting of \texttt{id}
and \texttt{measured\_at}. This still guarantees uniqueness while enabling
efficient time-based partitioning.

\begin{lstlisting}[language=PGTSSQL, caption={SQL commands to adjust the primary key of the "meas" table.}, label={lst:sql-create-hypertable}]
-- Drop existing primary key on meas
ALTER TABLE public.meas DROP CONSTRAINT IF EXISTS meas_pkey;

-- Add composite primary key (id, measured_at)
ALTER TABLE public.meas ADD CONSTRAINT meas_pkey PRIMARY KEY (id, measured_at);
\end{lstlisting}

After that, we created the hypertable on \texttt{meas} using the following
command:

\begin{lstlisting}[language=PGTSSQL, caption={SQL command to create a hypertable on the "meas" table.}, label={lst:sql-create-hypertable-convert}]
-- Convert meas to hypertable on measured_at
SELECT create_hypertable('public.meas', 'measured_at', chunk_time_interval => INTERVAL '1 day', if_not_exists => TRUE, migrate_data => TRUE);
\end{lstlisting}

In addition, we enabled compression on the \texttt{meas} table to reduce
storage requirements and improve query performance for older data:

\begin{lstlisting}[language=PGTSSQL, caption={SQL command to enable compression on the "meas" table.}, label={lst:sql-enable-compression}]
-- Enable compression
ALTER TABLE public.meas SET (
  timescaledb.compress,
  timescaledb.compress_segmentby = 'fk_user_id',
  timescaledb.compress_orderby = 'measured_at DESC'
);

-- Add compression policy (compress data older than 7 days)
SELECT add_compression_policy('public.meas', INTERVAL '7 days');
\end{lstlisting}

This enables compression for \texttt{meas}, using \texttt{fk\_user\_id} for
segmentation and ordering by \texttt{measured\_at} in descending order. The
compression policy then automatically compresses chunks older than seven days,
keeping recent data uncompressed for fast access while optimizing historical
storage and query performance.

\subsubsection{TimescaleDB in AirScout}

Overall, TimescaleDB proved to be an excellent fit for our project. Its
time-series optimizations, especially for IoT workloads, and its seamless
integration with PostgreSQL---a system we already knew well---made it a strong
choice. The benchmark results also indicate that TimescaleDB performs better
for our use case (see Table~\ref{tab:tsbs-query-results}). Combined with the
weighted decision matrix from Table~\ref{tab:airscout-decision-matrix}, this
further supported our decision to use it for AirScout.

In our project-specific integration tests, we inserted more than $400{,}000$
measurements into the database, which it handled without issue. In addition,
the separate TSBS benchmark discussed in Chapter~\ref{sec:benchmarking}
operated at a much larger scale. Based on both results, we are confident that
TimescaleDB will continue to scale as the project grows.